% Generated draft: requires mathematical and source review.
\documentclass[11pt]{article}
\usepackage[margin=25mm]{geometry}
\usepackage{fontspec}
\setmainfont{Latin Modern Roman}
\usepackage{amsmath,amssymb,mathtools,unicode-math}
\setmathfont{Latin Modern Math}
\usepackage{enumitem,needspace,xurl,hyperref}
\setlength{\parindent}{0pt}
\setlength{\parskip}{5pt}
\setlength{\emergencystretch}{3em}
\newcommand{\N}{\mathbb{N}}
\newcommand{\Z}{\mathbb{Z}}
\newcommand{\Q}{\mathbb{Q}}
\newcommand{\R}{\mathbb{R}}
\newcommand{\C}{\mathbb{C}}
\newcommand{\F}{\mathbb{F}}
\newcommand{\A}{\mathbb{A}}
\newcommand{\PP}{\mathbb{P}}
\newcommand{\E}{\mathbb{E}}
\newcommand{\eps}{\varepsilon}
\newcommand{\dd}{\,\mathrm{d}}
\begin{document}
\section*{UnsolvedMath ID: 30002260}
\textbf{Source problem code:} catalogue entry
\subsection{Definitions and mathematical statement}
The supplied record does not specify the ambient theory, the class of definable groups, or the permitted parameter sets. The following gives a conventional interpretation of its notation, conditional on these choices being intended.

Let \(T\) be a complete first-order theory, let \(\mathfrak C\models T\) be a sufficiently saturated and strongly homogeneous monster model with saturation cardinal \(\kappa\), and let \(A\subseteq\mathfrak C\) satisfy \(|A|<\kappa\). Let \(G^*\) be a group definable in \(\mathfrak C\), with its defining parameters included in \(A\). A subgroup has bounded index if its index in \(G^*\) is less than \(\kappa\). A set is type-definable over \(A\) if it is the set of realizations of a small collection of formulas with parameters from \(A\); it is \(A\)-invariant if every automorphism of \(\mathfrak C\) fixing \(A\) pointwise preserves it.

In the usual connected-component notation, \(G_A^{*00}\) is the smallest subgroup of \(G^*\) that is type-definable over \(A\) and has bounded index, and \(G_A^{*000}\) is the smallest \(A\)-invariant subgroup of bounded index. With these components interpreted in the usual monster-model setting, \(G_A^{*000}\) is normal in \(G_A^{*00}\), and the group quotient is
\[
Q(G^*,A)=G_A^{*00}/G_A^{*000}.
\]

\textbf{Classification question.} Characterize, up to abstract group isomorphism, all groups \(H\) for which there exist \(T,\mathfrak C,G^*,A\) as above such that
\[
H\cong Q(G^*,A).
\]
The unrestricted variation of these data and the use of abstract group isomorphism are explicit interpretive conventions: the supplied record does not establish whether narrower hypotheses or additional quotient structure are intended.

\textbf{Unresolved source qualification.} The original statement also asks whether such a quotient can be non-abelian. The supplied statement-verification note says that the same report contribution constructs non-abelian examples and that only the classification question remains open. In contrast, the supplied literature assessment dated 2026-08-21 says that non-abelianity remains undecided. These accounts conflict. The classification question is the target retained by the supplied corrected statement; the status of the original non-abelianity question requires source reconciliation.
\subsection{Short English statement}
Classify the groups that can arise as the quotient of a definable group's smallest parameter-type-definable subgroup of bounded index by its smallest parameter-invariant subgroup of bounded index. The ambient hypotheses and the intended notion of classification need confirmation. The supplied accounts disagree about whether the original non-abelianity question has already been resolved.
\subsection{Sources}
\begin{itemize}
\item Local registry: \texttt{data/unsolvedmath/problems.json}, permanent UnsolvedMath ID \(30002260\), source problem code \texttt{catalogue entry}, title \emph{Model-Theoretic Connected-Component Quotients}. The supplied record belongs to the Oberwolfach Reports Open Problems collection and gives proposed year 2013; no individual proposer is supplied.
\item \emph{Model Theory: Groups, Geometry, and Combinatorics} (2013), Oberwolfach Reports, pp.\ 5--66, DOI 10.4171/owr/2013/01. Supplied source citation and literature-source label: \emph{Model Theory, Oberwolfach Report 01/2013}. \url{https://doi.org/10.4171/owr/2013/01}.
\item \emph{Model-Theoretic Connected Components in Group Extensions}. Supplied literature reference: \url{https://arxiv.org/abs/1201.5221}.
\item Literature checked at: \texttt{2026-08-21}. This is supplied dated curation, not an independently verified assessment. It labels the problem open and states that connected-component work provides structural cases without classifying the quotients or deciding non-abelianity; the background expressly requests independent verification of time-sensitive status and cited claims.
\item The supplied record labels the statement \texttt{corrected\_verified}. Its verification note reports non-abelian examples in the same report contribution and retains only classification as open. That supplied verification claim conflicts with the dated literature assessment and has not been independently checked here.
\end{itemize}
\end{document}
